\documentclass{article}

\usepackage[margin=1in]{geometry}
\usepackage{fancyhdr}
\usepackage{graphicx}
\usepackage{color}

\definecolor{INCFBlue}{rgb}{0.0,0.59,1.0}

\definecolor{issuecolor}{rgb}{0.8,0.8,0.8}
\newcommand{\issue}[1]{%
\begin{center}
\colorbox{issuecolor}{\parbox{0.8\linewidth}{\textbf{Issue:} #1}}
\end{center}%
}

\newcommand{\suggestion}[2]{%
\begin{center}
\colorbox{issuecolor}{\parbox{0.8\linewidth}{\textbf{#1:} #2}}
\end{center}%
}

\begin{document}

\pagestyle{empty}

\begin{center}
\includegraphics[width=0.7\columnwidth]{./incf.png}
\end{center}

\vspace*{1cm}

\noindent\rule{\columnwidth}{1pt}
\noindent\rule{\columnwidth}{2pt}

\vspace*{1cm}

\begin{center}

\noindent{\Huge \bf Conceptual Summary:}\\

\vspace{0.5cm}
\noindent{\Large \bf Network Interchange format for NEuroscience (NineML): \\ User Layer}\\
\vspace{0.5cm}
\noindent{\large INCF Task Force on Multi-Scale Modeling}\\
\vspace{0.5cm}

\noindent{\large Version: 0.93}\\

\end{center}

\vspace*{0.75cm}

\noindent\rule{\columnwidth}{2pt}
\noindent\rule{\columnwidth}{1pt}

\vspace*{3cm}
\noindent{\Large

\begin{center}
{\bf Authors (in alphabetical order): }
\end{center}

\noindent Anatoli Gorchetchnikov \\

%\vspace*{0.5cm}

\noindent {\bf Date:} \today

}


\title{NineML (9ML) User Layer Summary}

\newpage
\pagestyle{plain}

\tableofcontents
\newpage

\abstract
This draft intends to describe our current opinions on User Layer concepts
of NineML.

This work currently summarizes the different ideas and propositions elaborated
during seven Task Force meetings and User Layer group discussions. It might
contain some misrepresentations or omissions of important issues and shall be
updated during future work.

\vskip 1in

\paragraph{Changes from version 0.92:}
\begin{itemize}
\item Adapted the format similar to the Abstraction layer specification.
I think this document shall essentially become user layer section of the
general specification document.
\item Attempted to reorganize and adapt the text to be more like ``Object
Model'' for the User Layer.
\end{itemize}

\paragraph{Changes from version 0.91:}
\begin{itemize}
\item Replaced ``node'' with ``component'' as was decided in Antwerp 2010
meeting.
\item Replaced references to ``core semantics'' with references to
``abstraction layer''
because the term ``core semantics'' is no longer used.
\item As Antwerp 2010 decided, properties are represented by a generic
tag and specified
by attribute {\tt name}.
\item Added the reference to units of measurements dimensionality checking
against
the abstraction layer as was decided in Antwerp 2010.
\item Moved edge effect descriptions from connectivity components to layout
components. Also
noted that some layouts might require complex masks.
\end{itemize}

\paragraph{Changes from version 0.9:}
\begin{itemize}
\item added footnotes for some issues raised in Stockholm June meeting
\item added brief description of interactions between user layer and
simulation software that does not support abstraction layer
\item removed random number generator nodes as decided in Stockholm
\item added space, region, and layout nodes as a first draft of geometrical
concepts (based on NetworkML proposal by Padraig and Robert as well as
Stockholm discussions)
\item removed recurrent projections that can be described on the population
level
\item cut out the previous approach to cell positioning within population
\end{itemize}
\newpage

\section{Introduction}

With an increasing number of studies related to large-scale neuronal
network modeling, the International Neuroinformatics Coordinating
Facility (INCF) has identified a need for standards and guidelines to
ease model sharing and facilitate the replication of results across
different simulators. To create such standards, the INCF has formed a
program on Multiscale Modeling to develop a common standardized
description language for neuronal network models.

The name of the proposed standard is NineML (Network Interchange for
Neuroscience Modeling Language) and its first version is aimed at
descriptions of large networks of spiking neurons.

\subsection{Language syntax}

It is envisioned that the language should not require a
specific syntax or rely on any given technology or platform.
Rather it is anticipated that the language can be
employed by defining the model elements in a variety of different
syntaxes including a native (domain specific) language, Python, Java,
and XML (for example).

\subsection{Scope of the Usel Layer}

The user layer provides the syntax for specifying the model and parameters to
be used to instantiate the key elements of a spiking neuronal network in terms
that are directly linked to the biological counterparts of the model. This
includes descriptions of individual elements (cells, synapses, inputs) and the
constructs for describing the grouping of these components into networks. In
addition the user layer defines the syntax for specifying a range of
connectivity patterns.

The remainder of this document summarizes the concepts that underlie user
layer design.

\section{Core Concepts}

User layer has eight key elements: data types, property, definition,
component, population, group, set, and projection. Out of these eight
only two -- component and group -- are used at the top level of the
description, descriptions of all the others are embedded in the higher
level element and are not used independently. The following subsections
discuss each of these elements.

\subsection{Data types}

Data types are the objects that allow the encapsulation of necessary data
and are not intended as standalone objects but rather as parts of composite
objects like {\tt property} or {\tt selector}.

\subsubsection{Dimension}

Dimension is an object that provides physical dimension like length,
mass, or time. A set of predefined dimensions is provided in the
NineML library and the user can derive integral dimensions from these
(for example velocity is derived from length and time). A special case
of a {\tt Dimension} is a predefined empty dimension for dimensionless
quantities.

\subsubsection{String}

String is an object that contains a string of text.

\issue{This goes well with Hugo's idea that a name shall be a property like
all others rather then an attribute. Maybe also for the notes.}

\subsubsection{URL}

URL is an object that contains a reference to an internet resource in a form
of a URL.

\issue{Use it for the notes and links to AL definitions?}

\subsubsection{Enumerated}

\issue{This seems needed for properties that can take values like
convergent-divergent, etc.}

\subsubsection{Boolean}

Boolean is an object that contains a binary true-false value or a logical
expression that can be evaluated into such a value.

\subsubsection{Quantity}

Quantity is a compound data type object that encapsulates a numerical value
and a unit of measurements. Unit has to be of one of the {\tt Dimension} type.
Any numeric quantity in the language has to be of this type,
dimensionless quantities shall use predefined empty {\tt Dimension}.

\begin{table}[htb]
\center
\begin{tabular}{|c|}
\hline
\hline
Quantity \\
\hline
\hline
{\em unit}: {\tt Dimension} \\
\hline
{\em value}: \{Number$|$Function$|$Random Distribution\} \\
\hline
\end{tabular}
\end{table}

There are two kinds of quantities: values of the first kind are given to the
model by the user and stay fixed, values of a second kind are computed within
the model during simulation. For all practical reasons the syntax of the user
layer descriptions is identical for both kinds of quantities. Furthermore,
because NineML does not provide any default values for quantities, it is a
job of the user to provide initial values for all defined quantities. To
ensure the integrity of the model NineML requires all initial values to be
set in the user layer description. For batch simulations and other
modifications any of the values given in the user layer can be overwritten
by a simulation setup description, but this is outside of the scope of the
current version of NineML.

Some quantities can have values drawn from random distribution. In this case
user layer description of the quantities includes a reference to a random
distribution component (section \ref{randDist}) instead of a numeric value.
Other quantities might be calculated according to some function dependent on
the geometry of the system or some other quantities. These can be defined by
including inline abstraction layer definitions or MathML. Note that quantities
that depend on geometry will remain uninstantiated until the network is
created.

\subsection{Property}

Property is a construct that combines a value of one of the data types
defined in the previous subsection, an indicator of the used type, and a
label. The label should match the corresponding label in the abstraction
layer definition.

\begin{table}[htb]
\center
\begin{tabular}{|c|}
\hline
\hline
Property \\
\hline
\hline
{\em label}: {\tt String} \\
\hline
{\em value}: \{{\tt Quantity}$|${\tt Boolean}$|${\tt Enumerated}\} \\
\hline
{\em note}: \{{\tt String}$|${\tt URL}\} (optional)\\
\hline
\end{tabular}
\end{table}

The user can set the value of the property and (when applicable for the data
type, e.g. {\tt Quantity}) the units of measurement. These units are also
checked against the dimensionality of the corresponding property definition
in the abstraction layer.

User can add a note to a property. These notes are
intended to provide a specific reference to the research paper page
where the component is described and similar kind of information. These
notes are not intended to duplicate the abstraction layer documentation
or any other documentation, thus they shall not provide mathematical
description and other details of the component implementation. Note can
contain text (type {\tt String}) or link to an Internet resource (type
{\tt URL}).

\subsection{Definition}

All constructs in the user layer have their mathematical or algorithmical
definitions in the abstraction layer. Definition is an object that
establishes a link between user layer and abstraction layer. The initial
version of NineML allows to put references to external (including user
space abstraction layer) definitions.

\begin{table}[htb]
\center
\begin{tabular}{|c|}
\hline
\hline
Definition \\
\hline
\hline
{\em language}: {\tt String} \\
\hline
{\em link}: {\tt URL}\\
\hline
\end{tabular}
\end{table}

The language should be flexible enough to allow representation of concepts
that do not yet exist, as it is developed to serve the forefront of research.
A simple mechanism to add concepts that are not part of the standard is
provided through external (other than abstraction layer based) definitions.
It is the choice of a simulator developer to support these definitions during
initial stage of NineML development. It is unclear at this point who will
eventually convert these definitions into native NineML definitions, but the
maturation of NineML shall eliminate the need to support simulator-specific
definitions.

\subsection{Reference}

Reference is an object that can replace {\tt Definition} object in situations
where user needs to reuse the same definition for multiple instances of
objects. In this case the first instance shall be described using
{\tt Definition} and the rest can use the first one through {\tt Reference}.
Label in the {\tt Reference} must match a label in a previously defined
object.

\begin{table}[htb]
\center
\begin{tabular}{|c|}
\hline
\hline
Reference \\
\hline
\hline
{\em label}: {\tt String} \\
\hline
\end{tabular}
\end{table}

\subsection{Component}

The basic building block of NineML is called component. In the user layer
description component is a reference to an object defined in the abstraction
layer. Abstraction layer defines the mathematics of the component and this
definition is then referred in the user layer through an object of type
{\tt Definition}.

\begin{table}[htb]
\center
\begin{tabular}{|c|}
\hline
\hline
Component \\
\hline
\hline
{\em label}: {\tt String} \\
\hline
{\em definition}: \{{\tt Definition}$|${\tt Reference}\}\\
\hline
{\em note}: \{{\tt String}$|${\tt URL}\} (optional)\\
\hline
\colorbox{issuecolor}{\parbox{0.4\linewidth}
{\center Set of {\tt Property} objects}} \\
\hline
\end{tabular}
\end{table}

In addition to the {\tt Definition} object, the component encapsulates
user-given ID or label, and a set of {\tt Property} objects. The composition
of this set of properties is defined in the abstraction layer (or externally)
and instantiated in the user layer description of component.

\issue{The documentation shall provide the mapping between mathematical
description of the object and the corresponding properties labels.}

User can add short notes to each component's description similar to
notes described above for properties.

As a result, a complete description of a component in the user layer consists
of three parts: reference to a definition, a list of properties, and an
optional note.

To reduce the size of the resulting network description user can refer to
already described component by user-given name instead of providing link to
abstraction layer or external definition. In this case the properties of the
component that have to be redefined are stated explicitly, the properties that
are inherited from the original description are omitted.

If the simulator only supports user layer, then the simulator developers
can create mappings directly between the reference to a definition in the user
layer description of the node and the intrinsic simulator code that implements
the same mathematics.

\subsubsection{Random Distribution Components}
\label{randDist}

This component allows to define random distribution with corresponding
parameters. These component allows the reuse of the same distribution
multiple times similar to reuse of all other components.

\subsubsection{Neuronal Components}
\label{neuron}

The description of the neuronal component defines a prototypical neuron
which can be reused multiple times within the network. As a consequence,
multiple connectivity patterns can be applied to this neuron, which
results in a different set of synaptic inputs. Therefore, no description
of these inputs shall be provided at the component level, they shall be
described at the projection level (Section \ref{projections}). Simulation
software shall take care of complete construction of each neuron by
analyzing both levels of description.

\subsubsection{Plasticity Components}
\label{plasticity}

Plasticity components handle the synaptic weight and its possible
modification. Note that synaptic weight is a dimensionless quantity and
as such always have units of the predefined empty {\tt Dimension}. This
is done to allow same plasticity rules to operate on
different types of postsynaptic responses. Synaptic weight defined in the
plasticity component determines only the magnitude of the response, the
shape and physical properties are defined in the corresponding post-synaptic
response node (section \ref{secSynapse}).

\subsubsection{Post-synaptic Response (PSR) Components}
\label{secSynapse}

Post-synaptic response components define the effect imposed on the
post-synaptic cell dynamics by triggering the synaptic input. This
definition includes only the shape of this effect and does not include the
exact magnitude of the effect. The magnitude is described separately through
plasticity rules and synaptic weights in section \ref{plasticity}. Note
though that the units of the effect are defined here and synaptic weights in
plasticity nodes are left dimensionless.

\subsubsection{Connectivity Components}

Connectivity components contain the definition of a connectivity pattern
between neurons. Connectivity components are based on the layouts of the
source and target groups. All connectivity components have a delay property
that sets the axonal delay for the corresponding set of projections.

\subsubsection{Space Components}

Space components define the coordinate systems used to set up the geometry
of the network. These components can be Cartesian, polar, and maybe some other
coordinate systems. More than one space components can be defined in the same
network, for example Cartesian 3D for cortical cells and polar 2D for retinal
space.

\issue{In case we need rotated, translated, or scaled versions of the same
coordinate system, shall we define a new space component?}

\issue{In case user combines two models how shall the space be handled?}

\subsubsection{Region Components}

Region components define a finite region given a coordinate system. A
reference to a corresponding coordinate system as well as boundaries of a
region constitute the properties of a region component.

\subsubsection{Layout Components}
\label{layout}

Layout components define how neurons or groups of neurons are positioned
either within a given region of space or on spaceless indexed grid. Both
discrete grids and continuous mappings can operate on the same region, so
it is possible to create a 3D grid of neurons for one population as well
randomly distributed set of neurons from other population within the same
region. Spaceless indexed grids in addition to providing neuronal indices
include descriptions of the edge handling. In general, a layout is simply
a mapping from indices to generalized positions.

\issue{Some layouts require non-trivial masks that we started to work on in
Antwerp 2010, but this work is not completed.}

\subsection{Population}

Population describes a collection of \emph{identical} elements. To allow
the arbitrary complexity of the model, these elements can be not only
components but also groups (Section \ref{grouping}). Populations refer to
a spatial layout as well as to a prototypical component or group by a user
given ID (label). Different populations can use the same prototype if the user
wants to define two populations of same neurons or circuits but with
different connectivity or with different positional distribution. Referring
to the same layout from different populations is more tricky, since if it
is a fixed layout then the positions of neurons in two populations will
coincide.

\begin{table}[htb]
\center
\begin{tabular}{|c|}
\hline
\hline
Population \\
\hline
\hline
{\em label}: {\tt String} \\
\hline
{\em prototype}: \{{\tt Component}$|${\tt Group}\}\\
\hline
{\em number}: {\tt Quantity} \\
\hline
{\em layout}: {\tt Component->Layout} \\
\hline
{\em note}: \{{\tt String}$|${\tt URL}\} (optional)\\
\hline
\end{tabular}
\end{table}

\issue{Shall we limit the {\tt Component} subtypes that can be used
for prototypes? If so, to which subtypes aside from neuronal subtype?}

\issue{In some layouts the number of elements follows from the layout
structure.}

\issue{Shall we represent the number which is always natural and
dimensionless through different data type rather than {\tt Quantity}?}

Populations do not include connectivity information and are not used
independently, but rather as a part of the parent group (Section
\ref{grouping}).


\subsection{Constructive Group}
\label{grouping}

The grouping of objects is of two major types: \emph{\bf constructive}
grouping used to {\bf create} objects and \emph{\bf access} grouping used
to define the subgroups of {\bf already created} objects for the purposes
of connectivity or monitoring. The latter type is called {\tt Set} and
described in the Section \ref{accessGroup}. The former one is called
{\tt Group} and is is a two-layered structure including of a set of
populations on the lower level and connectivity between these populations
together with access groups on the upper level. Both the number of
populations and the number of elements within population can be set to
one. Processing of a constructive group by a NineML-compliant
software shall lead to creation of instances for all elements, define the
spatial structure of the group, and provide internal connectivity if
necessary.

\begin{table}[htb]
\center
\begin{tabular}{|c|}
\hline
\hline
Group \\
\hline
\hline
{\em label}: {\tt String} \\
\hline
{\em note}: \{{\tt String}$|${\tt URL}\} (optional)\\
\hline
\colorbox{issuecolor}{\parbox{0.4\linewidth}
{\center Set of {\tt Set} objects}} \\
\hline
\colorbox{issuecolor}{\parbox{0.4\linewidth}
{\center Set of {\tt Projection} objects}} \\
\hline
\colorbox{issuecolor}{\parbox{0.4\linewidth}
{\center Set of {\tt Population} objects}} \\
\hline
\end{tabular}
\end{table}

Groups allow hierarchical build of larger groups from smaller groups. Rather
than having a deep hierarchy, however, the structure is kept flat by
referencing. Each Group contains a user-given ID or label. Groups do not
contain the description of the elements that are repeated, but only refers
to their original description by these IDs. The scoping rules for these names
are described in section \ref{scope}.

The rules of construction provided by Groups shall not be
executed immediately on occurrence of such a group within the description.
The process of construction shall proceed from top-most Group in the hierarchy
down to component groups. This allows the flexibility of defining one
group in the description and reusing it multiple times within the network.
This also allows the inclusion of developed networks in the larger models
as sub-components. The top-most Group in the hierarchy is not marked anywhere
in the description, it is rather selected by the user through the software
interface. This way the user can simulate individual components of the network
for debugging and tuning purposes without modifying the rest of the model.

A collection of projection entries within group can describe the internal
connectivity within this group according to the rules shown in section
\ref{projections}. In case this connectivity depends on some features
described outside of the group it has to be defined in the point of the
hierarchy from where both the feature and question and this group are
directly accessible by descending the tree.

\subsection{Set}
\label{accessGroup}

Set is an access group and it does not define how the cells or subnetworks
are created, but rather how they shall be selected for the purpose of
monitoring of the the activity or building the connectivity. To eliminate
Sets that fully replicate corresponding Groups any Group
declaration also implicitly defines one Set of all the elements it contains
plus one Set for each of the populations in the Group. Group user-given IDs
and population user-given IDs can be used anywhere Set IDs can.
Declaring a set does not add any objects or components to a model
specification: it just forms a new, possibly empty, set from those
components that have already been declared.

\begin{table}[htb]
\center
\begin{tabular}{|c|}
\hline
\hline
Set \\
\hline
\hline
{\em label}: {\tt String} \\
\hline
{\em selector}: {\tt Select} \\
\hline
{\em note}: \{{\tt String}$|${\tt URL}\} (optional)\\
\hline
\end{tabular}
\end{table}

{\tt Select} is not fully defined at this point, it is intended to use
logical operations on per-element basis and return a subset of all
elements for which the final result of these operations is true.

\subsubsection{Selection by a User-given ID}

Matching criteria for these selections are user-given IDs that can be found
while descending the tree within the parent group of the set description.
When all matches are exhausted, the union of each such binding is
constructed and returned.

\subsubsection{Selection by a Cell Property}

This selection only makes sense for cells that do have certain property
defined. If later the model will be changed and a different
neuronal component will be used as a prototype, this selectшщт will become
invalid. The user shall be notified of this inconsistency by a simulation
software and resolve the problem manually. Furthermore, if the property
in question has units of measurements attached to it, the selection shall
compare it to a quantity that also includes units.

Finally, some properties are variable over
the course of each simulation, and if these properties are allowed in for
selection that means that the set will be different on every simulation
step. For now NineML only allows sets based on the properties that are
preset during network construction and remain constant through the
simulation.

\subsubsection{Selection by Geometry}

This type of selection takes valid positions as arguments. The
compliant software shall ensure that both positions are described within the
same coordinate system, an error message shall be raised if this is not the
case. For topographic projections the position of the source within its group
(population) partially determines the position of the target within a
different group (population). In this case we might need more complicated
coordinate mapping then a simple reduction to a uniform coordinate space
(e.g. cells with certain 2D retinal position project to cells with certain
3D coordinates in the cortical sheet).

\subsubsection{Selection by Logical Combinations of Criteria}

Within each select statement logical operations are possible. For completeness
all three (AND, OR, NOT) are allowed.

\subsection{Projection}
\label{projections}

A projection holds a description of the connectivity between two sets
of cells. Similar to the set, the projection does not create any new
components in the network. The purpose of projection object is to bind
source and destination using a certain connectivity scheme, plasticity
rule and post-synaptic response. Projection is not used standalone,
only within a group. Along with a user-given ID or name unique inside
the hosting group, each projection description contains references to
source and target sets. Each projection description leads to creation
of a set of projections rather than a single projection.

\begin{table}[htb]
\center
\begin{tabular}{|c|}
\hline
\hline
Projection \\
\hline
\hline
{\em label}: {\tt String} \\
\hline
{\em source}: \{{\tt Set}$|${\tt Group}\} \\
\hline
{\em destination}: \{{\tt Set}$|${\tt Group}\} \\
\hline
{\em synapse}: {\tt Component->PSR} \\
\hline
{\em plasticity}: {\tt Component->Plasticity} \\
\hline
{\em connectivity}: {\tt Component->Connectivity} \\
\hline
{\em note}: \{{\tt String}$|${\tt URL}\} (optional)\\
\hline
\end{tabular}
\end{table}

Projection description also includes references to a plasticity component
that controls the synaptic weight (section \ref{plasticity}) and post-synaptic
response node that controls the influence of the input through this projection
on the post-synaptic cell dynamics (section \ref{secSynapse}). Both appear in
the description of the projection rather than neuron because the same type of
neuron can use different synapses when instantiated in different populations,
similarly the same postsynaptic response can be used in multiple projections
with different plasticity rules. Instead of providing the user with all
possible pre-wired combinations, NineML allows user to combine the
plasticity$\rightarrow$response$\rightarrow$neuron chain from a set of small
standard components.

Finally, a connectivity pattern or rule shall be set for a projection through
reference to a connectivity component.

\section{Structure of a User Layer Description}

In order to simplify the descriptions themselves and the mechanisms for
combining multiple components of the model the user layer description
is consisting of multiple files that define various components and
contains the syntax to import external files.

\subsection{Scope of User Layer Descriptions}
\label{scope}

In order to avoid the collisions of user-given IDs as well as to
provide the user with the ability to redefine some of the imported
components the following set of rules is suggested:
\begin{itemize}
\item{Every identifier is considered unique only within its parent}
\item{To access the identifier a full path from the point of access
to this identifier shall be provided.}
\item{Any imported construct can be redefined within the importing
description by adding a new definition of the same object. All properties
of the imported object are imported from external file and the values of
some (or even all) of them are overwritten by the local description.}
\item{If multiple redefinitions of the same entity are found, then the most
recent one will be used. The most recent is determined by the order of
parsing of the top level model description file. Here the rule shall be
enforced that when multiple files are included all nodes are parsed before
all groups. This will allow a smooth transition from networks with simple
nodes and groups to networks where some nodes and some groups are redefined.}
\item{If several different subfields are redefined in different places of the
description, then all of these new definitions will be used. Please note that
this is not the optimal way, and the better way to achieve the same effect is
to provide all redefinitions in a unique location. Simulation software
developers can (but at this point of time are not obliged to) optimize the
NineML code that is output by their software by grouping all redefinitions
of components of the same node into a single node.}
\end{itemize}

\newpage
\appendix

\part*{Appendix}
\addcontentsline{toc}{part}{Appendix}

\section{INCF Task Force}

An oversight committee has been setup to guide the direction and
objectives of the initiative with membership as follows:

\begin{itemize}

    \item Chair: Erik De Schutter, Okinawa Institute of Science and Technology, Japan
    \item Ad Aertsen, Bernstein Center, Freiburg, Germany
    \item Upinder Bhalla, National Center for Biological Sciences, Bangalore, India
    \item Nicolas Brunel, CNRS-Université René Descartes, Paris, France
    \item Sharon Crook, Arizona State University, Tempe, USA
    \item Gustavo Deco, Universitat Pompeu Fabra, Barcelona, Spain
    \item Markus Diesmann, RIKEN, Wako, Japan
    \item Tomoki Fukai, RIKEN, Wako, Japan
    \item Wulfram Gerstner, LCN, EPFL, Lausanne, Switzerland
    \item Stephen Grossberg, Boston University, MA, USA
    \item Rolf Kötter, Radboud University, Nijmegen, The Netherlands
    \item Anders Lansner, Royal Institute of Technology (KTH), Stockholm, Sweden
    \item Peter Latham, Gatsby Computational Neuroscience Unit, London, UK
    \item Christiane Linster, Cornell University, Ithaca, NY, USA
    \item Wolfgang Maass, Technische Universität Graz, Austria
    \item Ernst Niebur, Johns Hopkins University, Baltimore, MD, USA
    \item Felix Schürmann, EPFL, Lausanne, Switzerland
    \item Simon Thorpe, Faculté de Médecine de Rangueil, Toulouse, France
    \item Alessandro Treves, International School for Advanced Studies (SISSA), Trieste, Italy
    \item Shiro Usui, RIKEN, Wako, Japan
    \item Mark van Rossum, University of Edinburgh, UK
    \item Xiao-Jing Wang, Yale University, New Haven, CT, USA

\end{itemize}


A task force has been assembled to evaluate and implement the
recommendations of the oversight committee with membership as follows:

\begin{itemize}

    \item Lead: Sean Hill, BMI, EPFL, Lausanne, Switzerland
    \item Lead: Anatoli Gorchetchnikov, Boston University, USA
    \item Lead: Ivan Raikov, Okinawa Institute of Science and Technology, Japan
    \item Board Secretary: Yann Le Franc, University of Antwerp, Belgium
    \item Robert Cannon, Textensor Limited, Edinburgh, UK
    \item Robert Clewley, Georgia State University, Atlanta, USA
    \item Hugo Cornelis, University of Texas Health Science Center, San Antonio, USA
    \item Andrew Davison, CNRS, Gif-sur-Yvette, France
    \item Erik De Schutter, Okinawa Institute of Science and Technology, Okinawa, Japan
    \item Mikael Djurfeldt, INCF Secretariat and PDC,
Royal Institute of Technology, Stockholm, Sweden
    \item Padraig Gleeson, University College London, UK
    \item Michael Hines, Yale University, New Haven, CT, USA
    \item Birgit Kriener, Max Planck Institute for Dynamics and Self-Organization,
Göttingen, Germany

    \item Chung-Chuan Lo, National Tsing Hua University, Hsinchu, Taiwan
    \item Abigail Morrison, Bernstein Center Freiburg, Germany
    \item Eilif Muller, LCN, EPFL, Lausanne, Switzerland
    \item Hans Ekkehard Plesser, Norwegian University of Life Sciences, Aas, Norway
    \item Subhasis Ray, National Center for Biological Sciences, Bangalore, India
    \item Lars Schwabe, University of Rostock, Germany
    \item Botond Szatmary, Neurosciences Institute, San Diego, CA, USA

\end{itemize}

and thus includes the authors or primary contributors to a wide
variety of computational neuroscience efforts including the Blue Brain
Project, GENESIS-3, KInNeSS, MOOSE, NEURON, NEST, PyNN and NeuroML.

\bibliographystyle{apalike}
\bibliography{CondensedUL}

\end{document}

The user layer of NineML is intended to be primarily machine-readable and uses XML
syntax. It is designed with a focus on ease of parsing, verification, and
automatic model construction. This does not prevent advanced users from editing
the user layer XML descriptions by hand, but the primary means for creation of
these descriptions is expected to be software tools that will convert GUI- or
script-based representations of objects and properties into valid XML.

For specific examples of implementing these concepts please refer to
a General NineML Description.

\subsection{Basics of the Structure, Layouts, and Coordinates}

Declaring the existence of a population in a constructive group implicitly
introduces a notion of indexing for its elements, so, for example selection rules
may use conditions on the indexes to select subsets of cells. Population
contains a reference to a layout node to map indices onto coordinates,
which can also be used in selection rules for access groups.

